\chapter{Технологический раздел}

\section{Выбор языка и среды программирования}

В качестве языка программирования был выбран язык \texttt{C}~\cite{cdocs}, так как он является основным языком для разработки драйверов ядра Linux и обеспечивает прямой доступ к низкоуровневым функциям системы.

Для компиляции драйвера используется \texttt{gcc}~\cite{gccdocs} с флагами, соответствующими версии ядра Linux. В качестве среды разработки выбран текстовый редактор \texttt{Neovim}~\cite{nvimdocs} с плагинами для работы с кодом ядра.

Сборка драйвера осуществляется с помощью \texttt{Makefile}, который автоматизирует процесс компиляции и учитывает зависимости от версии ядра.

\section{Функциональность разработанного ПО}

Разработанный драйвер преобразует контроллер Xbox Series S в устройство ввода типа "мышь" с расширенными возможностями. Основная функциональность включает:

\begin{enumerate}
	\item Назначение кнопок контроллера на функции мыши:
	\begin{itemize}[label=---]
		\item Кнопка A — левая кнопка мыши;
		\item Кнопка В — правая кнопка мыши;
		\item Кнопка X - средняя кнопка мыши.
	\end{itemize}
	\item Функцию перемещения курсора мыши с помощью стиков геймпада.
	\item Регулируемые параметры чувствительности:
	\begin{itemize}[label=---]
		\item Чувствительность движения курсора;
		\item Чувствительность прокрутки;
		\item Размер мертвой зоны;
	\end{itemize}
	\item Автоматическое определение контроллера при подключении через USB.
\end{enumerate}

Драйвер работает в фоновом режиме, автоматически определяя подключение контроллера и обеспечивая бесшовную интеграцию с существующей системой ввода Linux.

\section{Структура данных драйвера}

Основные структуры данных драйвера:

\begin{lstlisting}[label=lst:shandler,caption=Структура драйвера Xbox Series S, language=c]
struct usb_xpad {
    struct input_dev *dev;          // Основное устройство ввода (геймпад)
    struct usb_device *udev;        // USB-устройство
    struct usb_interface *intf;     // USB-интерфейс
    
    struct xpad_mouse *mouse;       // Структура для эмуляции мыши
    unsigned char *idata;           // Буфер входящих данных
    unsigned char *odata;           // Буфер исходящих данных
    
    int mapping;                    // Карта преобразования ввода
    int xtype;                      // Тип контроллера
    const char *name;               // Имя устройства
};

struct xpad_mouse {
    struct input_dev *dev;          // Устройство мыши
    bool left_button;               // Состояние левой кнопки
    bool right_button;              // Состояние правой кнопки  
    bool middle_button;             // Состояние средней кнопки
    int deadzone;                   // Мертвая зона для стиков
};
\end{lstlisting}

\section{Инициализация USB-драйвера}

\begin{lstlisting}[label=lst:usbinit,caption=Инициализация USB-драйвера, language=c]
static struct usb_device_id xpad_table[] = {
    { USB_DEVICE(0x045e, 0x0b12) }, // Xbox Series S|X Controller
    {}
};

static struct usb_driver xpad_driver = {
    .name = "xpad_xbox_series",
    .probe = xpad_probe,
    .disconnect = xpad_disconnect,
    .suspend = xpad_suspend,
    .resume = xpad_resume,
    .id_table = xpad_table,
};
\end{lstlisting}

Поля структуры \texttt{usb\_driver}:
\begin{itemize}[label=---]
	\item \texttt{name} — имя драйвера, отображаемое в системных логах;
	\item \texttt{probe} — функция, вызываемая при обнаружении поддерживаемого устройства;
	\item \texttt{disconnect} — функция, вызываемая при отключении устройства;
	\item \texttt{suspend/resume} — функции управления энергосбережением;
	\item \texttt{id\_table} — таблица идентификаторов USB-устройств.
\end{itemize}

\section{Параметры модуля}

Драйвер поддерживает настройку через параметры модуля:
\begin{itemize}[label=---]
	\item \texttt{mouse\_emulation} — включение/выключение эмуляции мыши;
	\item \texttt{mouse\_sensitivity} — чувствительность движения курсора (1-10);
	\item \texttt{scroll\_sensitivity} — чувствительность прокрутки (1-10).
\end{itemize}

\section{Листинг загружаемого модуля ядра}

Исходный код загружаемого модуля и файла сборки представлены в приложениях А и Б соответственно.